% Portada de J3.3: las cuatro factorizaciones como productos de bloques.
\documentclass[tikz,border=0pt]{standalone}
\usepackage{resonancia}
\begin{document}
\begin{tikzpicture}
  \clip[rounded corners=6pt] (0,0) rectangle (12,3.6);
  \shade[top color=rAzul!16, bottom color=white] (0,0) rectangle (12,3.6);
  \tikzset{blk/.style={draw=#1!70, fill=#1!22, line width=0.7pt, minimum width=0.85cm, minimum height=0.85cm}}
  % A = L U
  \node[blk=rAzul] at (0.6,2.4) {};  \node at (1.25,2.4) {$=$};
  \node[blk=rVerde] at (1.9,2.4) {}; \node[blk=rNaranja] at (2.85,2.4) {};
  \fill[rVerde!70] (1.475,2.825) -- (1.475,1.975) -- (2.325,1.975) -- cycle;
  \fill[rNaranja!70] (2.425,2.825) -- (3.275,2.825) -- (3.275,1.975) -- cycle;
  \node[etiqueta] at (1.7,1.45) {\small $\mathbf A=\mathbf L\mathbf U$};
  % A = Q R
  \node[blk=rAzul] at (4.9,2.4) {};  \node at (5.55,2.4) {$=$};
  \node[blk=rVioleta] at (6.2,2.4) {}; \node[blk=rNaranja] at (7.15,2.4) {};
  \fill[rNaranja!70] (6.725,2.825) -- (7.575,2.825) -- (7.575,1.975) -- cycle;
  \node[etiqueta] at (6.0,1.45) {\small $\mathbf A=\mathbf Q\mathbf R$};
  % A = U S V'
  \node[blk=rAzul] at (9.0,2.4) {};  \node at (9.65,2.4) {$=$};
  \node[blk=rVioleta] at (10.3,2.4) {}; \node[blk=rVerde, minimum width=0.85cm] at (11.2,2.4) {};
  \draw[color=rTinta, line width=1.1pt] (10.775,2.825) -- (11.625,1.975);
  \node[etiqueta] at (10.2,1.45) {\small $\mathbf A=\mathbf U\boldsymbol\Sigma\mathbf V^{\mathsf T}$};
  \node[etiqueta, color=rGris] at (6.0,0.6) {\small LU · Cholesky · QR · SVD: una matriz, bloques más sencillos};
\end{tikzpicture}
\end{document}
